\scanpage{44}

\uline{Revêtements quadratiques}

\diagram{\begin{tikzcd}[ampersand replacement=\&,column sep=small]0\arrow[r]\&\underline O\arrow[r]\& A\arrow[r,"\varphi"]\&\underset{\mathrm{inv}}{\check L}\arrow[r]\&0\end{tikzcd}\quad $E=\check A$\quad\begin{tikzcd}[ampersand replacement=\&,column sep=small]0\arrow[r,dash]\&L\arrow[r]\&E\arrow[r,"\psi"]\&\underline O\arrow[r]\&0\end{tikzcd}}

\[
\begin{aligned}
\underline P&=\psi^{-1}(1)\quad\text{tors. sous }L &\psi(x)&=\langle1,x\rangle\\
\underline P_\lambda&=\psi^{-1}(\lambda)\quad\text{id}&&=x(1)
\end{aligned}
\]

$\bullet\quad T\in\Gamma\underline P_2\subset E$\qquad fonction trace

$\bullet\quad N\in\Gamma\operatorname{Sym}^2(E)\subset\operatorname{Sym}^{\unsure{\bullet|*}}(E)$\qquad forme quadratique \uline{norme} (fonctions pol. sur $A$)

NB\qquad 1)

\diagram{\begin{tikzcd}[ampersand replacement=\&]0\arrow[r]\&E\otimes L\arrow[r]\&\operatorname{Sym}^2(E)\arrow[r,"\chi"]\&\underline O\arrow[r]\&0\end{tikzcd}}

\[
\begin{aligned}
&\text{où }\chi(Q)=Q(1)\\
&\text{on a }N\in\chi^{-1}(1)\quad(\text{tors. sous }E\otimes L)\\
&\text{i.e. }N(1)=1
\end{aligned}
\]

\deleted{\unsure{b)}} 2) \deleted{\unsure{On}} On a

\diagram{\begin{tikzcd}[ampersand replacement=\&]0\arrow[r,dash]\& L^{\otimes2}\arrow[r]\&E\otimes L\arrow[r]\&L\arrow[r]\&0\end{tikzcd}}

et $\operatorname{Sym}^2(E)$ a suite de comp. à quotients successifs $\underline O$, $L$, $L^{\otimes2}$. On identifie $L^{\otimes2}$ à ss-mod de $\operatorname{Sym}^2(E)$

\margin{attention au signe $\longrightarrow$ $\delta$ est le discr. de $N$ \uline{changé de signe}}

$\bullet\quad-\delta\in\Gamma(L^{\otimes2})$\qquad ``discriminant'' de $N$

\uline{NB}\qquad Forme bil. sym $\varphi_N$ associée à $N$

\[
\begin{aligned}
\varphi_N:\quad A\times A&\mathbin{-}\underline O\qquad\text{soit}\qquad A\longrightarrow\check A\overset{\mathrm{déf}}{\simeq}E\\
\text{d'où}\qquad \underset{\substack{\wr\\\check L}}{\bigwedge^2 A}&\longrightarrow\underset{\substack{\wr\\L}}{\bigwedge^2 E}\qquad\text{d'où }\delta\in\Gamma\underbrace{L^{\otimes2}}_{\substack{\wr\\\underline{\operatorname{Hom}}(\check L,L)}}
\end{aligned}
\]

% slanted
\margin{\uline{NB} On a $\varphi_N(x,y)=x\bar y+y\bar x$, et $\operatorname{Tr}(x)\operatorname{Tr}(y)=\varphi_N(x,y)+\operatorname{Tr}(xy)$ \unsure{ou i.e.|on trouve} $\varphi_N(x,y)=\operatorname{Tr}(x)\operatorname{Tr}(y)-\operatorname{Tr}(xy)$.}

\[
\boxed{T^2-4N=\inserted{+}\delta}\quad\text{i.e}\quad\boxed{4N=T^2\deleted{\unsure{+}}{\unsure{-}}\delta}.
\]

\margin{$\uparrow$ Attention au signe\,!}

\uline{NB} Soient $x',x''\in\Gamma E$ deux sections \uline{distinctes} en tt pt, et qui correspondent à des homs d'alg. $A\overset{x',x''}{\rightrightarrows}\underline O_X$ (donc $x',x''\in\Gamma\underline P_1$). On a donc $A\simeq\underline O_X\times\underline O_X$

\[
\begin{aligned}
\deleted{\unsure{x'(x+y)|x'(a,b)}}\qquad x'(a,b)&=a,\qquad x''(a,b)=b\\
\operatorname{Tr}(a,b)&=a+b,\qquad N(a,b)=ab
\end{aligned}
\]

d'où

\[
\boxed{\begin{aligned}T&=x'+x''\\N&=x'x''\\\delta&=(x'-x'')^2\end{aligned}}\qquad\deleted{\unsure{\delta=(x'-x'')^2}}
\]

% archive sheet number: 43
